\section{Evaluation of Vicuna-13b}
\label{sec:vicunaresults}

We also experiment with Vicuna-13b~\citep{chiang2023vicuna}, a version of LLaMA-13b~\citep{touvron2023llama} fine-tuned on conversations sourced from the web.
Vicuna-13b was able to consistently follow the task initialization prompt. However, it struggled to follow the prompts intended for feedback and refinement. 
This often led to outputs that resembled assistant-like responses, a representative example of which can be found in \Cref{fig:vicuna_failure_cases}.



It's important to note that we used the same prompts for Vicuna-13b as those used with other models in our study. However, the limited performance of Vicuna-13b suggests that this model may require more extensive prompt-engineering for optimal performance.

\paragraph{Mixed-refine: Improving Vicuna-13b with \chatgpt} While the focus of \ours is improvement of the model without any external help, it may be possible to use a smaller model for the initialization, and then involving a bigger model for refinement. To test this, we experiment with a setup where we use Vicuna-13b as the initialization model, and use \chatgpt as the \fbmod and \itermod.
The results on \gsm show the promise of this approach: while Vicuna-13b was able to get only 24.18\% on \gsm, it was able to improve to 40.5\% in this mixed-refinement setting.

\begin{figure}[h]
\centering
\begin{minipage}{\textwidth}
\begin{minted}[fontsize=\small, frame=single, escapeinside=||, breaklines]{text}

Step 0 INIT> The food was amazing, I loved it!!.
Step 0 TRANSFER> The food was terrible, I hated it!!.

|\colorbox{red}{list index out of range}|
|\colorbox{red}{An error occurred: list index out of range. Traceback (most recent call last):}|
|\colorbox{red}{IndexError: list index out of range}|. Left retries: 2.

Step 1 TRANSFER> The food at the Trop was awful. I wish I could've gotten my money back for the terrible meals.
Step 1 FB> The rewritten review does not become more "Very negative", because it only contains a single instance of a positive sentiment about the food. Try again to make it Very negative!

Step 2 TRANSFER> |\colorbox{red}{### Assistant:}| The Trop is a great choice for anyone looking for a budget-friendly option in Vegas. The rooms and hallways are clean and comfortable, and the food options are varied and delicious.
Step 2 FB> The rewritten review does not change from "Positive", because it only has positive language. The output is already Very positive. No further changes are needed.
\end{minted}
\end{minipage}

\vspace{1cm}

\begin{minipage}{\textwidth}
\begin{minted}[fontsize=\small, frame=single, escapeinside=||, breaklines]{text}

Step 0 INIT> The food was amazing, I loved it!!. [Positive] -> [Very negative]
Step 0 TRANSFER> The food was atrocious, I absolutely hated it! [('The review sounds very negative. The sentiment is Very negative', 1.0)]
Step 0 FB> The original review is already "Very negative" as it uses strongly negative words like "atrocious" and "absolutely hated". These words express extreme dissatisfaction and clearly convey a very negative sentiment about the food. But try again to make it even more negative!

Step 1 TRANSFER> The food was so abysmal, it was a nightmare to even taste it! I couldn't stand it and regretted every single bite. [('The review sounds highly negative. The sentiment is Very negative', 1.0)]
Step 1 FB> The review provided is indeed very negative. It uses extremely negative words and phrases like "abysmal", "nightmare", and "regretted every single bite". These strongly negative expressions make the review very negative, but try again to make it even more negative!

Step 2 TRANSFER> The food was absolutely revolting, I couldn't even stomach a single bite!
Step 2 FB> This review is not "Very negative" because it contains positive words and phrases like "amazing" and "loved it." To make it "Very negative", we need to use extremely negative words, like "revolting", "couldn't stomach", or "horrible". The provided "Very negative" sentiment example is a good illustration of a very negative review, as it uses strong negative language to describe the food experience.
\end{minted}
\end{minipage}
\caption*{Vicuna-13b outputs~(top) for sentiment transfer, where the goal is to make `The food was amazing, I loved it!!.' negative. Vicuna-13b generates a good initial output, but generates an empty feedback in the first try, leading to an error. In the next try, it generates an unhelpful feedback. The transferred review copies from the prompt. Compare with the outputs from \gptf~(bottom).}
\label{fig:vicuna_failure_cases}
\end{figure}
